\documentclass{article}
\usepackage[utf8]{inputenc}
\usepackage[T1]{fontenc}
\usepackage{geometry}
\usepackage{graphicx}
\usepackage{float}
\usepackage{hyperref}
\usepackage[english,indonesian]{babel}

\title{Optimasi Pemrosesan Template LaTeX Menggunakan Pendekatan Berbasis AI}
\author{Academic Support Agent}
\date{\today}

\begin{document}
\maketitle

\begin{abstract}

\end{abstract}

\chapter{PENDAHULUAN}


\chapter{TINJAUAN PUSTAKA}


% Original content (commented out):
% \documentclass[a4paper,10pt]{article}

% Paquetes
\usepackage{fontspec}
\usepackage{titlesec}
\usepackage{titling}
\usepackage{setspace}
\usepackage[hidelinks]{hyperref}
\usepackage{graphicx}
\usepackage{caption}
\usepackage{subcaption}
\usepackage{float}
\usepackage{xcolor}
\usepackage[affil-it]{authblk}
\usepackage{biblatex}
\addbibresource{Mendeley.bib}

% Idioma
\usepackage[spanish]{babel}
\usepackage[utf8]{inputenc}


% Configuraciones
\setmainfont{Times New Roman}
\hypersetup{
    colorlinks,
    linkcolor={red!50!black},
    citecolor={blue!50!black},
    urlcolor={blue!80!black}
}


\begin{document}
\begin{titlepage}
\centering
\begin{figure}[t]
	\centering
	\includegraphics[scale=0.15]{utn.jpg}
    \vspace{0.5cm}
\end{figure}%
	{\LARGE Proyecto Final\par}
    {\LARGE 2016\par}
	\vspace{1cm}
	{\huge\bfseries Titulo del Proyecto\par}
	\vspace{1cm}
    {\LARGE Abstract\par}
    \vspace{1cm}
	{\Large\itshape Integrante 1, Integrante 2, Integrante 3\par}
    \vspace{2cm}
\begin{flushleft}
	\textbf{Palabras clave: incluir al menos 5 palabras clave separadas por coma.}
\end{flushleft}
	\vfill
\end{titlepage}

\section{Desarrollo}
\label{sec:desarrollo}
Un abstract debe ser un breve resumen de su proyecto. Si está bien escrito, el mismo debe motivar al lector a saber más sobre su trabajo.
\paragraph{}
Existen componentes claves que no pueden dejar de figurar en cualquier abstract, que son:
\begin{itemize}
	\item \textbf{Motivación o problema:} ¿Qué es lo que motiva el desarrollo del proyecto? ¿Qué necesidad viene a satisfacer?
    \item \textbf{Métodos:} ¿Cómo planean resolverlo? ¿Qué enfoque utilizarán? Si bien en esta etapa aún no se encuentran definidos de manera final estos asuntos, deben al menos plantear cómo pretenden a priori resolver el problema.
    \item \textbf{Resultados:} ¿Qué resultados planean obtener? ¿A dónde pretenden llegar desarrollando el producto o proyecto de investigación?
    \item \textbf{Conclusiones:} Esta sección normalmente se completa una vez efectuado el trabajo. Aquí deben desarrollarse los hallazgos más importantes del proyecto de investigación o desarrollo del producto.
\end{itemize}

Es importante destacar que cada una de estas secciones son más o menos importantes dependiendo si se trata de un proyecto de investigación o del desarrollo de un producto y varía también dependiendo a qué área de aplicación se oriente cada uno.
\paragraph{}
Cabe destacar que no es necesario que los ítems sean desarrollados en forma de lista. Puede escribirse un texto que cubra todos los aspectos.
\paragraph{}
A diferencia de un paper, donde el abstract debe contener cerca de 200 palabras dentro de un mismo párrafo, en este caso la extensión deberá ser de una carilla aproximadamente y puede estar separado en varios párrafos. \cite{MendeleySupportTeam2011MendeleyGuide}

\newpage
\section{Tutores externos}
\label{sec:tutores}
Deben brindar una lista de los tutores externos que colaborarán con ustedes en el proyecto, indicando un breve resumen de cada uno.

\begin{itemize}
	\item Tutor 1: información...
    \item Tutor 2: información...
\end{itemize}

%Bibliografia

%\begin{thebibliography}{9}
%\bibitem{bibref}
%	Nombre Deautor,
%    Ciudad, Pais,
%    \textit{"Título"},
%    Año.
%\end{thebibliography}
\printbibliography
\end{document}

\end{document}